\chapter{Requirements Analysis}
\label{ch:requirements}

% BUDGET: 2 pages.

Before any code was written, a requirements analysis questionnaire was administered to
academic staff. Its purpose was not to confirm that the proposed system was a good idea, but
to establish what practising researchers would actually require of such a tool --- and, more
usefully, what would stop them using it. Several of its findings changed the design
substantially, and one of them remains the project's most significant unmet requirement.

\section{Method}

Participants were recruited from academic staff at the University of Liverpool. The study was
restricted to lecturers, professors and other PhD holders on the grounds that supervisory
experience is what makes a respondent able to judge a tool of this kind: a supervisor sees
both their own research workflow and the failure modes of a doctoral researcher learning it.
Recruitment was by direct invitation, and participation was voluntary and unpaid.

The instrument was an online questionnaire of five sections --- background, current workflow,
reactions to the proposed system, prospective use cases, and desired features --- combining
five-point Likert items with free-text questions, and taking roughly 10--15 minutes. It opened
with an information page describing the project and the intended use of responses, and
required explicit consent before any question could be answered. Respondents could skip any
question and stop at any time. No names, email addresses or other identifying fields were
collected. The full instrument is reproduced in Appendix~\ref{app:ra-instrument} and the
anonymised responses in Appendix~\ref{app:ra-responses}.

Two members of academic staff completed the questionnaire, in July~2026, before implementation
began. They are referred to below as Respondent~A and Respondent~B, and their disciplines are
generalised to \emph{mathematical sciences} because the combination of exact field, position
and institution would otherwise be identifying within a small population.

\begin{itemize}[nosep]
  \item \textbf{Respondent~A}: Senior Lecturer / Associate Professor, 6--10 years post-PhD,
        currently supervises doctoral researchers, uses LLM tools \emph{often} and across
        every part of their work.
  \item \textbf{Respondent~B}: University Teacher, 2--5 years post-PhD, does not supervise,
        and \emph{never} uses LLM tools.
\end{itemize}

Two respondents is a small sample, and no statistical claim is made from it. Its value lies
elsewhere: the two occupy opposite ends of the disposition towards this technology, so where
they agreed, the agreement is informative, and where they disagreed, the disagreement exposed
a design requirement that neither alone would have revealed.

\section{Findings}

Table~\ref{tab:ra-likert} gives the Likert responses on the proposed system.

\begin{table}[htbp]
\centering
\small
\caption{Responses to the attitudinal items (1 = strongly disagree, 5 = strongly agree).}
\label{tab:ra-likert}
\begin{tabularx}{\textwidth}{Xcc}
\toprule
\textbf{Statement} & \textbf{A} & \textbf{B} \\
\midrule
A tool like this could save significant time                              & 3 & 4 \\
I would trust a literature review from such a system without checking it  & \textbf{1} & \textbf{1} \\
I would be comfortable letting it run and modify experiment code autonomously & 4 & 1 \\
I would want to review and approve outputs at every stage                 & 4 & 5 \\
I am concerned about originality overlapping with existing work           & 4 & 5 \\
I would be comfortable using its outputs with appropriate review          & 4 & 3 \\
\midrule
Importance of flagging overlap with existing literature (1--5)            & \textbf{5} & \textbf{5} \\
\bottomrule
\end{tabularx}
\end{table}

\paragraph{Nobody will trust unchecked output.} The only item on which both respondents chose
an extreme of the scale, and the same extreme, was trusting a generated literature review
without checking it: both strongly disagreed. Both also wanted to review and approve outputs at
every stage. This is the single most consequential finding in the study, and it reframes the
design problem. A system that asks to be trusted has already failed; the requirement is that
its outputs be \emph{checkable}, and cheaply so.

\paragraph{Verifiability matters more than autonomy.} Respondent~B, who does not use LLM tools
at all, located the objection precisely: \emph{``By far, it's the interpretability of the LLM.
I want to know why the LLM made a specific decision, I cannot accept a black box when it comes
to research.''} Respondent~B also declined to endorse the premise of the tool in general
terms, remarking that it is \emph{``good for lazy people''} and closing with a direct
challenge: \emph{``I'm very cynical about LLMs in research, and your project can either cater
to my concerns, or ignore me and focus on other customers.''} That challenge was taken
seriously, and Section~\ref{sec:ra-decisions} describes what was done about it.

\paragraph{The two respondents wanted incompatible things.} Respondent~A judged all stages
equally valuable to automate, was comfortable with autonomous code execution, and wanted a
checkpoint at every one of the five stages. Respondent~B wanted exactly one stage automated
--- code generation and experiment execution --- was strongly opposed to autonomous code
execution, wanted a checkpoint only at report writing, and stated plainly: \emph{``I would
consider a tool like this for programming, but not academic writing.''} Notably, the stage
Respondent~B was willing to delegate is the one their free-text answers identified as their
own bottleneck (\emph{``my own programming abilities''}), while the stages they refused to
delegate are those they regard as constitutive of scholarship.

\paragraph{The multi-agent premise was itself questioned.} Respondent~A challenged the role
metaphor inherited from Agent Laboratory~\cite{schmidgall2025agentlab} directly: \emph{``I
don't know whether having PhD student as a role particularly works, as in their PhD they need
to do every stage of the process\ldots{} I don't know whether you need individual agents for
all the tasks anyway, can't you just have one LLM that does it all?''}

\paragraph{Both wanted the same things visible during a run.} Both respondents selected every
available option: live progress through each stage, intermediate outputs, confidence or
uncertainty estimates, source citations with links to the retrieved papers, and compute cost
and time estimates. On interface, Respondent~B preferred a simple web interface and
Respondent~A was indifferent --- against a proposal that had specified a command-line tool
with a graphical interface only if time permitted.

\section{Requirements derived, and what they changed}
\label{sec:ra-decisions}

Six design decisions follow directly from the above, and are traceable to it.

\begin{description}[style=unboxed,leftmargin=0pt,itemsep=0.35em]

\item[D1 --- Personas were replaced by functional decomposition.] The proposal specified
agents named for academic roles (PhD Student, ML Engineer, Postdoc, Reviewer). Respondent~A's
objection is correct: the metaphor describes a team structure that does not match how doctoral
research is actually done, and it says nothing about what each component is responsible for.
The implemented system therefore names its seven agents after the function each performs ---
Literature, Interdisciplinary Literature, Hypothesis, Experiment Planner, Coder, Writer,
Reviewer --- and defines each by an explicit input/output contract rather than by a persona.

\item[D2 --- Separate agents were retained, for a reason the questionnaire supplied.]
Respondent~A's follow-up question (``can't you just have one LLM that does it all?'') deserves
a direct answer, and the answer is a consequence of the trust finding. Stage boundaries are
where deterministic validation can be placed: because each agent's output is a validated
contract, the pipeline can check citations against retrieved papers, verdicts against executed
code, and coverage against the hypothesis set. A single undifferentiated model has nowhere to
put those checks.

\item[D3 --- Verification was designed in preference to trust.] Given that neither respondent
would accept unchecked output, and that Respondent~B's objection was specifically about
opacity, every decision the system makes that \emph{can} be settled mechanically is settled in
code rather than by the model: citation markers resolve only against papers actually retrieved,
hypothesis verdicts are computed from execution results, and results computed on synthesised
inputs withhold a verdict entirely. Chapter~\ref{ch:design} develops this as the central design
principle of the system.

\item[D4 --- The web interface was promoted from optional to core.] Given the interface
preference and the demand for stage-by-stage visibility, the graphical interface ceased to be
a stretch goal and became the primary means of interaction, showing live per-stage progress,
intermediate outputs, retrieved papers, and every draft as it is produced.

\item[D5 --- Co-pilot control became stage selection, not merely stage pausing.] The
respondents' incompatible preferences cannot be satisfied by a single pause-and-confirm
behaviour, because they disagree about \emph{which} stages should run at all. The implemented
system therefore lets a user run a contiguous range of stages rather than the whole pipeline,
supply their own papers instead of searching for them, and choose which hypothesis is carried
forward. Respondent~B's stated preference --- experiment code only, no automated writing ---
is expressible as a configuration rather than requiring a different tool.

\item[D6 --- Overlap detection was specified, and not delivered.] Both respondents rated the
importance of flagging overlap with existing literature at 5 out of 5, the only item to score
the maximum from both, and both were independently concerned about originality. This
corroborates Gupta and Pruthi's finding~\cite{gupta2025glitters} that a substantial share of
agent-generated research documents contain undetected plagiarism. It is reported here rather
than in the design chapter because it was \emph{not} implemented; it is the highest-priority
unmet requirement in the project, and Section~\ref{sec:future-overlap} returns to it.

\end{description}

Two further requests were also not met: confidence or uncertainty estimates, and compute cost
and time estimates, both of which every respondent wanted visible during a run. Neither is
present in the delivered interface, and both are discussed in Chapter~\ref{ch:conclusion}.

Finally, both respondents indicated willingness to evaluate a working prototype, and the
evaluation reported in Chapter~\ref{ch:evaluation} draws on that undertaking.
